% GENERATED FILE. Edit canonical lessons, then run npm run generate:books.
% canonical-source-hash: fnv1a64:01652c976c6985d9
% canonical-lessons: TA-C84-words, TA-C84-lines, TA-C84-mudhal-vaasippu, TA-R84-letters-that-look-alike, TA-R84-vowels-signs-and-the-dot, TA-R84-words-the-script-pages-built, TA-R84-the-u-family-completed

\chapter{Reading — Words, Lines, and a First Passage}
\label{ch:ta-mudhal-vaasippu}
\hlchaptermodality{84}

\begin{chapteropening}
\textbf{By the end of this chapter:} \emph{I can read a Tamil word as shapes rather than sounds, recognise a question by its last word, and hold a whole passage without translating it word by word.}

Thirty words of Tamil with no new word in them, four of whose lines have no verb at all because Tamil needs no word for is.
\end{chapteropening}

\section[(reading)]{(reading) — six words, counted in shapes}

\label{lesson:TA-C84-words}



\begin{quote}
Every word below you have written. Not one is new.

What is new is that nobody says them first.
\end{quote}

\begin{tcolorbox}[breakable,skin=enhanced,colback=orange!6,colframe=orange!45!black,boxrule=0.5pt,arc=1mm,left=6pt,right=6pt,top=4pt,bottom=4pt,fonttitle=\bfseries,title={Read this}]
\begin{quote}
\ta{வீடு} \ta{புத்தகம்} \ta{தண்ணீர்} \ta{பால்} \ta{அம்மா} \ta{அப்பா}
\end{quote}

Read down the list once, without stopping.

Again — and this time count the shapes in each word rather than the sounds.
\end{tcolorbox}

\subsection*{What to know first}
They do not match. Tamil joins a consonant and its vowel into one shape, so \textbf{\ta{வீடு}} is four sounds and two shapes.

That is the thing a reader has to stop fighting. You learned these words as sounds and wrote them as shapes, and reading asks you to take the shape whole rather than to unpack it back into sounds.

\textbf{\ta{அம்மா}} and \textbf{\ta{அப்பா}} are the clearest pair: same length, same rhythm, one shape different in the middle.

\subsection*{Your turn}
\begin{itemize}
\raggedright
  \item \emph{Say it:} the six words down the list, once, without stopping
\end{itemize}

\emph{Twice through:}

\begin{tcolorbox}[breakable,colback=teal!4,colframe=teal!35!black,title={Before you move on}]
How many of the six were new? (\textbf{None.})

Does a Tamil word have as many shapes as sounds? (\textbf{No} — fewer.)
\end{tcolorbox}
\section[(reading)]{(reading) — six lines, and a question that ends where it asks}

\label{lesson:TA-C84-lines}



\begin{quote}
The last lesson gave you words. These are things you would say.

You have said every one of them.
\end{quote}

\begin{tcolorbox}[breakable,skin=enhanced,colback=orange!6,colframe=orange!45!black,boxrule=0.5pt,arc=1mm,left=6pt,right=6pt,top=4pt,bottom=4pt,fonttitle=\bfseries,title={Read this}]
\begin{quote}
\ta{வணக்கம்}. \ta{நன்றி}. \ta{ஆம்}. \ta{இல்லை}. \ta{உன்} \ta{பெயர்} \ta{என்ன}? \ta{இங்கே}.
\end{quote}

Read down once, without stopping.

Again — and look at where the question word sits in the fifth line.
\end{tcolorbox}

\subsection*{What to know first}
At the end. \textbf{\ta{உன்} \ta{பெயர்} \ta{என்ன}?} — \emph{your name what?} English puts the question word first and Tamil puts it last, so a reader does not know a line is a question until the line is over.

That is worth meeting on a page, where you can look again. In speech the intonation gives it away early; in writing only the final word does.

\textbf{\ta{வணக்கம்}} is a bowing, from the verb to bend — the same idea as namaste and namaskar, except Tamil kept the noun.

\subsection*{Your turn}
\begin{itemize}
\raggedright
  \item \emph{Say it:} the six lines, once, in order, without stopping
\end{itemize}

\emph{Twice through:}

\begin{tcolorbox}[breakable,colback=teal!4,colframe=teal!35!black,title={Before you move on}]
When do you find out a Tamil line is a question? (\textbf{At its last word.})

What is \textbf{\ta{வணக்கம்}}, literally? (\textbf{A bowing.})
\end{tcolorbox}
\section[(reading)]{(reading) — the first passage}

\label{lesson:TA-C84-mudhal-vaasippu}



\begin{quote}
Six words, then six lines. Now twelve of them together.

Thirty words, and not one of them is new.
\end{quote}

\begin{tcolorbox}[breakable,skin=enhanced,colback=orange!6,colframe=orange!45!black,boxrule=0.5pt,arc=1mm,left=6pt,right=6pt,top=4pt,bottom=4pt,fonttitle=\bfseries,title={Read this}]
\begin{quote}
\ta{வணக்கம்}. \ta{நான்} \ta{நல்லா} \ta{இருக்கிறேன்}. \ta{இது} \ta{என்} \ta{வீடு}. \ta{வீடு} \ta{பெரிய}. \ta{நான்} \ta{புத்தகம்} \ta{படிக்கிறேன்}. \ta{புத்தகம்} \ta{நல்ல}, \ta{ஆனால்} \ta{பெரிய}. \ta{அம்மா} \ta{இங்கே}. \ta{அப்பா} \ta{அங்கே}. \ta{தண்ணீர்}, \ta{பால்}, \ta{தேநீர்}. \ta{இன்று} \ta{நான்} \ta{படிக்கிறேன்}. \ta{உன்} \ta{பெயர்} \ta{என்ன}? \ta{நன்றி}.
\end{quote}

Read it once without stopping. Do not translate as you go — let the sentences arrive.

Now read it again, and notice which lines have no verb at all.
\end{tcolorbox}

\subsection*{What to know first}
Four of them. \textbf{\ta{வீடு} \ta{பெரிய}.} \textbf{\ta{அம்மா} \ta{இங்கே}.} Tamil does not need a word for \emph{is}, so a noun and what you are saying about it stand side by side and the sentence is finished.

A reader has to accept that a line can end without anything that looks like a verb. English cannot do it; Tamil does it constantly.

The one joiner here is \textbf{\ta{ஆனால்}}. The lesson that taught it also taught \textbf{\ta{ஆனா}}, the same word as people say it — this passage sits on the written side of that split, which is where reading lives.

This is the first whole passage the book has asked you to hold at once, and it works because everything in it was paid for.

\subsection*{Your turn}
\begin{itemize}
\raggedright
  \item \emph{Say it:} the passage aloud, once, without stopping
\end{itemize}

\emph{Twice through:}

\begin{tcolorbox}[breakable,colback=teal!4,colframe=teal!35!black,title={Before you move on}]
How many words in that passage were new? (\textbf{None.})

What does Tamil not need in \textbf{\ta{அம்மா} \ta{இங்கே}}? (\textbf{A verb.})
\end{tcolorbox}
\section[pannireṇṭu meyyeḻuttu]{(twelve consonants, four confusable sets) — sorted by ear, not by date}

\label{lesson:TA-R84-letters-that-look-alike}



\begin{quote}
Read \textbf{\ta{வ}}, then \textbf{\ta{க}}, and name each.
\end{quote}

\subsection*{Script: the three that are all n}
\noindent
\begin{tabularx}{\linewidth}{@{}>{\raggedright\arraybackslash}X>{\raggedright\arraybackslash}X>{\raggedright\arraybackslash}X@{}}
\toprule
\textbf{letter} & \textbf{sound} & \textbf{where the tongue goes} \\
\midrule
\textbf{\ta{ந}} & \emph{na} & teeth — the soft one \\
\textbf{\ta{ண}} & \emph{ṇa} & curled back, on the roof \\
\textbf{\ta{ன}} & \emph{ṉa} & the ridge behind the teeth \\
\bottomrule
\end{tabularx}

They arrived one per chapter, pages apart, and no page has put them side by side. That is the whole reason this set is hard: \textbf{Tamil hears three different letters where English hears one sound.} The difference is in where the tongue lands, and the shapes carry no hint of it.

\subsection*{Script: two l-letters and two r-letters}
\noindent
\begin{tabularx}{\linewidth}{@{}>{\raggedright\arraybackslash}X>{\raggedright\arraybackslash}X@{}}
\toprule
\textbf{pair} & \textbf{letters} \\
\midrule
the l-sounds & \textbf{\ta{ல}} \emph{la}, ordinary · \textbf{\ta{ழ}} the curled one Tamil is named for \\
the r-sounds & \textbf{\ta{ர}} \emph{ra}, a tap · \textbf{\ta{ற}} \emph{ṟa}, harder and further back \\
\bottomrule
\end{tabularx}

\textbf{\ta{ழ}} deserves the pause. It is the sound in the language's own name, it has no English equivalent, and a reader who met it once, thirty chapters ago, will not place it cold now.

\subsection*{Script: five that stand alone}
No confusable twin, no set to belong to: \textbf{\ta{வ}} \emph{va} · \textbf{\ta{ப}} \emph{pa} · \textbf{\ta{ம}} \emph{ma} · \textbf{\ta{க}} \emph{ka} · \textbf{\ta{ச}} \emph{ca}.

Each one carries a built-in \emph{a}, as every bare Tamil consonant does. Reading one aloud means saying two things — the consonant and the vowel inside it.

\subsection*{Your turn}
Cover the tables.

\begin{itemize}
\raggedright
  \item \emph{Read it:} \textbf{\ta{ந}}, \textbf{\ta{ண}}, \textbf{\ta{ன}} — and say where the tongue goes for each
  \item \emph{Read it:} \textbf{\ta{ல}}, then \textbf{\ta{ழ}} — and say which one names the language
  \item \emph{Read it:} \textbf{\ta{ர}}, then \textbf{\ta{ற}} — and say which is the harder one
  \item \emph{Read it:} \textbf{\ta{வ}}, \textbf{\ta{ப}}, \textbf{\ta{ம}}, \textbf{\ta{க}}, \textbf{\ta{ச}} — five in a row, no hesitating
  \item \emph{Say it:} the vowel hiding inside a bare Tamil consonant
\end{itemize}

\begin{tcolorbox}[breakable,colback=teal!4,colframe=teal!35!black,title={Before you move on}]
How many n-letters does Tamil write, and which of them curls the tongue back? (\textbf{Three; \ta{ண}}.) Which letter gives the language its name?
\end{tcolorbox}
\section[uyireḻuttum kuṟiyum puḷḷiyum]{(three vowels, five signs, one dot) — the same vowel, standing or hanging}

\label{lesson:TA-R84-vowels-signs-and-the-dot}



\begin{quote}
Read \textbf{\ta{அ}}. Then look at \textbf{\ta{்}} and say what it does to the letter it sits on.
\end{quote}

\subsection*{Script: standing, or hanging}
\noindent
\begin{tabularx}{\linewidth}{@{}>{\raggedright\arraybackslash}X>{\raggedright\arraybackslash}X>{\raggedright\arraybackslash}X@{}}
\toprule
\textbf{vowel} & \textbf{on its own} & \textbf{hanging on a consonant} \\
\midrule
long \emph{ā} & \textbf{\ta{ஆ}} & \textbf{\ta{ா}} \\
\emph{u} & \textbf{\ta{உ}} & \textbf{\ta{ு}} \\
short \emph{a} & \textbf{\ta{அ}} & nothing — it is already inside \\
\bottomrule
\end{tabularx}

That table is the whole system. A Tamil vowel gets a \textbf{full-size letter} when nothing comes before it, and a \textbf{hanging sign} when a consonant does. Short \emph{a} is the exception that proves it: every bare consonant already contains one, so it needs no sign at all.

\subsection*{Script: the three signs with no standing partner here}
\textbf{\ta{ி}} \emph{i} · \textbf{\ta{ே}} \emph{ē} · \textbf{\ta{ோ}} \emph{ō} — met one per chapter, and each one replaces the built-in \emph{a} of whatever letter it lands on.

Read them as a set and the shapes stop being arbitrary: \textbf{\ta{ே}} and \textbf{\ta{ோ}} share the same upward stroke, because \textbf{\ta{ோ}} is that stroke with a tail. Learning one buys most of the other.

\subsection*{Script: the dot that removes a vowel}
Three consonants were never gathered with the rest: \textbf{\ta{ட}} \emph{ṭa} · \textbf{\ta{த}} \emph{ta} · \textbf{\ta{ய}} \emph{ya}. Each carries its built-in \emph{a}, and \textbf{\ta{்}} takes it away.

So one mark and one letter give you two readings — \textbf{\ta{த}} is \emph{ta}, \textbf{\ta{த்}} is a bare \emph{t}. Everything the signs above do is a swap; this is the only one that subtracts.

\subsection*{Your turn}
Cover the tables.

\begin{itemize}
\raggedright
  \item \emph{Read it:} \textbf{\ta{ஆ}} and \textbf{\ta{ா}} — and say what the difference between them is
  \item \emph{Read it:} \textbf{\ta{உ}} and \textbf{\ta{ு}} — the same question
  \item \emph{Read it:} \textbf{\ta{ி}}, \textbf{\ta{ே}}, \textbf{\ta{ோ}} — name each
  \item \emph{Read it:} \textbf{\ta{ட}}, \textbf{\ta{த}}, \textbf{\ta{ய}} — and say the vowel inside each
  \item \emph{Say it:} what \textbf{\ta{்}} does
  \item \emph{Write it:} \textbf{\ta{வ}}, the first shape this book ever asked you to copy
\end{itemize}

\begin{tcolorbox}[breakable,colback=teal!4,colframe=teal!35!black,title={Before you move on}]
When does a Tamil vowel take a full-size letter rather than a hanging sign? (\textbf{When no consonant comes before it.}) Which mark subtracts a vowel instead of swapping one?
\end{tcolorbox}
\section[vācikka kaṟṟa ēḻu col]{(seven words you learned to read before you learned to spell them)}

\label{lesson:TA-R84-words-the-script-pages-built}



\begin{quote}
Read \textbf{\ta{சரி}}, then \textbf{\ta{என்}}. You have said both since the first two chapters.
\end{quote}

\begin{tcolorbox}[breakable,skin=enhanced,colback=orange!6,colframe=orange!45!black,boxrule=0.5pt,arc=1mm,left=6pt,right=6pt,top=4pt,bottom=4pt,fonttitle=\bfseries,title={Read this}]
\noindent
\begin{tabularx}{\linewidth}{@{}>{\raggedright\arraybackslash}X>{\raggedright\arraybackslash}X>{\raggedright\arraybackslash}X@{}}
\toprule
\textbf{word} & \textbf{what it means} & \textbf{what it was really teaching} \\
\midrule
\textbf{\ta{இல்லை}} & no, there is not & the \textbf{ai} sign, which stands to the LEFT \\
\textbf{\ta{சரி}} & fine, all right & \textbf{\ta{ச}}, one letter for a whole family of sounds \\
\textbf{\ta{என்}} & my & a word-initial vowel, reached by a rule you had \\
\textbf{\ta{பெயர்}} & name & a second sign that stands to the left of its letter \\
\textbf{\ta{எப்படி}} & how & the retroflex \textbf{\ta{ட}} \\
\textbf{\ta{இருக்கிறீர்கள்}} & are you (respectful) & the \textbf{\ta{ு}} sign, in a word of eight pieces \\
\textbf{\ta{தமிழ்}} & Tamil & \textbf{\ta{ழ}}, and the name of the language \\
\bottomrule
\end{tabularx}

Every one was a word you could already say. The writing pages never asked you to decode something unfamiliar — they put a known word in front of you and let the shapes attach to a sound you owned.
\end{tcolorbox}

\subsection*{Script: the signs that sit on the wrong side}
Two of the seven carry the same surprise. In \textbf{\ta{இல்லை}} and in \textbf{\ta{பெயர்}}, a vowel sign is written \textbf{before} the consonant it is pronounced after.

That is worth holding because it breaks the habit every reader brings: the order on the page is not the order in the mouth. A Tamil reader takes in the group and then says it.

\subsection*{Script: \ta{த} and \ta{ழ}}
\textbf{\ta{தமிழ்}} is two letters worth stopping on. \textbf{\ta{த}} is the soft dental \emph{t}, the one made against the teeth. \textbf{\ta{ழ}} is the curled sound with no English partner, and the language wears it in its own name.

Reading the name correctly is the single clearest proof that the script has stopped being decoration.

\subsection*{Your turn}
Cover the table and work from the left column only.

\begin{itemize}
\raggedright
  \item \emph{Read it:} \textbf{\ta{இல்லை}}, \textbf{\ta{சரி}}, \textbf{\ta{என்}} — three in a row
  \item \emph{Read it:} \textbf{\ta{பெயர்}}, then \textbf{\ta{எப்படி}}
  \item \emph{Read it:} \textbf{\ta{இருக்கிறீர்கள்}} — slowly, group by group
  \item \emph{Read it:} \textbf{\ta{தமிழ்}}, and name its last letter
  \item \emph{Say it:} which two of the seven write a vowel sign to the left of its letter
\end{itemize}

\begin{tcolorbox}[breakable,colback=teal!4,colframe=teal!35!black,title={Before you move on}]
Which letter gives Tamil its name, and which side of its consonant does the \textbf{ai} sign sit on? (\textbf{\ta{ழ}}; the left.) Which of the seven is the longest on the page?
\end{tcolorbox}
\section[u-kuṭumpam muḻumai]{(\ta{உ} \ta{ஊ} \ta{ு} \ta{ூ} — and the two that open o)}

\label{lesson:TA-R84-the-u-family-completed}



\begin{quote}
Read \textbf{\ta{உணவு}}, then \textbf{\ta{ஊர்}}. Say what the first letter of each is doing differently.
\end{quote}

\subsection*{Script: the square the u-family always formed}
Four pieces, arriving in four separate chapters, and they were a grid the whole time:

\noindent
\begin{tabularx}{\linewidth}{@{}>{\raggedright\arraybackslash}X>{\raggedright\arraybackslash}X>{\raggedright\arraybackslash}X@{}}
\toprule
\textbf{} & \textbf{standing alone} & \textbf{hanging on a consonant} \\
\midrule
\textbf{short u} & \textbf{\ta{உ}} & \textbf{\ta{ு}} \\
\textbf{long ū} & \textbf{\ta{ஊ}} & \textbf{\ta{ூ}} \\
\bottomrule
\end{tabularx}

\textbf{\ta{உணவு}} holds both columns of the top row — it opens with the standing \textbf{\ta{உ}} and closes with the hanging \textbf{\ta{ு}}, the same vowel wearing two shapes in one word. That word is the argument for the grid, and it is the reason this page exists: no chapter could show the square, because no chapter had all four corners.

\begin{tcolorbox}[breakable,skin=enhanced,colback=orange!6,colframe=orange!45!black,boxrule=0.5pt,arc=1mm,left=6pt,right=6pt,top=4pt,bottom=4pt,fonttitle=\bfseries,title={Read this}]
\noindent
\begin{tabularx}{\linewidth}{@{}>{\raggedright\arraybackslash}X>{\raggedright\arraybackslash}X>{\raggedright\arraybackslash}X@{}}
\toprule
\textbf{word} & \textbf{means} & \textbf{the piece it brought} \\
\midrule
\textbf{\ta{உணவு}} & food & standing \textbf{\ta{உ}} against hanging \textbf{\ta{ு}} \\
\textbf{\ta{ஊர்}} & town, home place & standing \textbf{\ta{ஊ}}, the long partner \\
\textbf{\ta{மூன்று}} & three & the hanging \textbf{\ta{ூ}}, closing the family \\
\textbf{\ta{பேசு}} & speak & the long-e sign \textbf{\ta{ே}} \\
\bottomrule
\end{tabularx}
\end{tcolorbox}

\subsection*{Script: o, standing and hanging}
\textbf{\ta{ஒன்று}} (\textquotedblleft{}one\textquotedblright{}) opens with \textbf{\ta{ஒ}}, the standing short o. \textbf{\ta{போ}} (\textquotedblleft{}go\textquotedblright{}) carries \textbf{\ta{ோ}}, the hanging long one — and \textbf{\ta{ோ}} is the odd shape of the set, because it is built from \textbf{two} marks, one placed before its consonant and one after.

So a single vowel can be written on both sides of the letter it belongs to at once. Reading it means taking the whole group in before saying anything.

\subsection*{Your turn}
Cover the grid.

\begin{itemize}
\raggedright
  \item \emph{Say it:} the four corners of the u-square, and which corner each shape is
  \item \emph{Read it:} \textbf{\ta{உணவு}} — and point to both u's in it
  \item \emph{Read it:} \textbf{\ta{ஊர்}}, \textbf{\ta{மூன்று}}, \textbf{\ta{பேசு}}
  \item \emph{Read it:} \textbf{\ta{ஒன்று}}, then \textbf{\ta{போ}}
  \item \emph{Say it:} which sign is built from two marks, and where each mark sits
\end{itemize}

\begin{tcolorbox}[breakable,colback=teal!4,colframe=teal!35!black,title={Before you move on}]
Which word on this page shows the same vowel standing and hanging at once? (\textbf{\ta{உணவு}}.) And which sign wraps around its consonant on both sides?
\end{tcolorbox}
